\chapter{Carbonylverbindingen: nucleofiele addities, acetalen en bescherming}\label{ch:b1:acetals-protection}

Los glucose op in water en bijna niets ervan blijft over als het
openketenaldehyde dat in leerboeken getekend wordt: het molecuul buigt
om, zijn eigen hydroxylgroep addeert aan zijn aldehydegroep, en er sluit
zich een zesring. Dezelfde additie van een alcohol aan een carbonylgroep,
nog eens herhaald met een zure \omterm{def:b1:elementary-steps:catalyst}{katalysator}, geeft \omterm{def:b1:acetals-protection:hemiacetal}{acetalen}, verbindingen
die bestand zijn tegen basen en tegen de reactiefste \omterm{def:b1:electronic-effects:nucleophile}{nucleofielen}. Die
stabiliteit is een werktuig: een aldehyde dat door een \omterm{def:b1:organomagnesium:organometallic}{grignardreagens}
vernietigd zou worden, kan als \omterm{def:b1:acetals-protection:hemiacetal}{acetaal} verborgen worden, door de reactie
meegevoerd, en op het eind teruggewonnen. Dit hoofdstuk bestudeert de
carbonylgroep als \omterm{def:b1:electronic-effects:nucleophile}{elektrofiel}, de addities van water en alcoholen, en de
strategie van de \omterm{def:b1:acetals-protection:protecting-group}{bescherming}.

\begin{recall}
Grignardaddities aan carbonylverbindingen (\cref{ch:b1:organomagnesium});
het \omterm{def:b1:extent-q-and-k:quotient}{reactiequotiënt} $Q$ en de evolutie van een systeem zolang $Q \neq K$
(\cref{ch:b1:extent-q-and-k}); \omterm{def:b1:electronic-effects:curly-arrow}{gebogen pijlen}, \omterm{def:b1:electronic-effects:nucleophile}{nucleofielen} en
\omterm{def:b1:electronic-effects:nucleophile}{elektrofielen} (\cref{ch:b1:electronic-effects}). Het schoolboek (jaar
12) introduceerde het idee een groep tijdens een synthese te beschermen.
\end{recall}

\section{De carbonylgroep als elektrofiel}

De C=O-binding is naar de zuurstof toe gepolariseerd
(\cref{ch:b1:periodicity}), en een \omterm{def:b1:lewis-resonance-vsepr:resonance}{resonantiestructuur} \ce{C+-O-} zet een
volledige positieve lading op het koolstofatoom: het koolstofatoom is
\omterm{def:b1:electronic-effects:nucleophile}{elektrofiel} en wordt door \omterm{def:b1:electronic-effects:nucleophile}{nucleofielen} loodrecht op het vlak van de groep
aangevallen. Sterke \omterm{def:b1:electronic-effects:nucleophile}{nucleofielen} (\omterm{def:b1:organomagnesium:organometallic}{grignardreagentia}, hydride, cyanide)
adderen rechtstreeks. Zwakke \omterm{def:b1:electronic-effects:nucleophile}{nucleofielen} --- water, alcoholen --- hebben
hulp nodig.

\begin{definition}[Elektrofiele activering]\label{def:b1:acetals-protection:activation}
De \emph{elektrofiele activering}\index{elektrofiele activering} van een
carbonylgroep is haar protonering (of haar coördinatie aan een \omterm{def:b1:lewis-resonance-vsepr:lewis-acid}{lewiszuur})
op de zuurstof. Het kation \ce{C=OH+}, waarvan de positieve lading via de
\omterm{def:b1:lewis-resonance-vsepr:resonance}{resonantiestructuur} \ce{C+-OH} met het koolstofatoom gedeeld wordt, wordt
zelfs door zwakke \omterm{def:b1:electronic-effects:nucleophile}{nucleofielen} aangevallen.
\end{definition}

\begin{omfigure}
\setchemfig{atom sep=2.1em}
\schemestart
\chemfig{C(-[:150]R)(-[:210]H)=O}
\arrow{->[\ce{H+}]}
\chemfig{C(-[:150]R)(-[:210]H)=O^{+}-H}
\arrow{<->}
\chemfig{C^{+}(-[:150]R)(-[:210]H)-O-H}
\schemestop

{\small Activering van een aldehyde door zuur: protonering van de
zuurstof levert een kation op waarvan een \omterm{def:b1:lewis-resonance-vsepr:resonance}{resonantiestructuur} de lading
op het koolstofatoom draagt.}
\end{omfigure}

\section{Hydraten en hemiacetalen}

Water addeert omkeerbaar aan aldehyden en ketonen, \ce{R2C=O + H2O <=>
R2C(OH)2}; het hydraat is bij ketonen meestal in de minderheid en kan
bij methanal overheersen. Een alcohol addeert op dezelfde manier.

\begin{definition}[Hemiacetaal en acetaal]\label{def:b1:acetals-protection:hemiacetal}
Een \emph{hemiacetaal}\index{hemiacetaal} is een verbinding waarin één
koolstofatoom zowel een OH- als een OR-groep draagt, gevormd door de
additie van een alcohol aan een carbonylgroep. Een
\emph{acetaal}\index{acetaal} is een verbinding waarin één koolstofatoom
twee OR-groepen draagt, gevormd uit een hemiacetaal en een tweede
alcoholmolecuul, met verlies van water.
\end{definition}

Openketenhemiacetalen zijn in hun evenwicht met het aldehyde en de
alcohol meestal in de minderheid; wanneer de OH en de C=O tot hetzelfde
molecuul behoren en een vijf- of zesring kan sluiten, overheerst het
cyclische \omterm{def:b1:acetals-protection:hemiacetal}{hemiacetaal}. Dat is het geval bij glucose, en bij het
eenvoudigere 5-hydroxypentanal.

\begin{omfigure}
\setchemfig{atom sep=2em}
\schemestart
\chemfig{HO-[:30]-[:-30]-[:30]-[:-30]-[:30]C(=[:90]O)-[:-30]H}
\arrow{<=>}
\chemfig{*6(O-(-[:-30]OH)-----)}
\schemestop

\medskip

{\small 5-Hydroxypentanal en zijn cyclische \omterm{def:b1:acetals-protection:hemiacetal}{hemiacetaal}, een zesring die
het zuurstofatoom bevat. De OH van koolstofatoom 5 heeft geaddeerd aan
het aldehyde van koolstofatoom 1; de ringvorm overheerst, zoals bij
glucose.}
\end{omfigure}

\section{Acetalen}

\begin{proposition}[Mechanisme van de acetaalvorming]\label{prop:b1:acetals-protection:acetal-mechanism}
In zuur milieu geven een aldehyde of keton en een alcohol het
\omterm{def:b1:acetals-protection:hemiacetal}{hemiacetaal} en daarna het \omterm{def:b1:acetals-protection:hemiacetal}{acetaal}, via de stappen: protonering van C=O;
additie van de alcohol; verlies van een proton (\omterm{def:b1:acetals-protection:hemiacetal}{hemiacetaal}); protonering
van de OH; verlies van water, wat een \omterm{def:b1:electronic-effects:intermediates}{carbokation} geeft dat door de
overblijvende OR-groep gestabiliseerd wordt (een oxocarbeniumion);
additie van een tweede alcohol; verlies van een proton. Elke stap is
omkeerbaar.
\end{proposition}

\begin{proof}
Elke stap is een zuur-baseoverdracht van een proton of de additie (of
het verlies) van een \omterm{def:b1:electronic-effects:nucleophile}{nucleofiel} aan (of van) een \omterm{def:b1:electronic-effects:nucleophile}{elektrofiel}
koolstofatoom, zoals beschreven in \cref{ch:b1:electronic-effects}. Het
verlies van water is mogelijk omdat het gevormde kation,
$\mathrm{R_2C^+{-}OR}$, een \omterm{def:b1:lewis-resonance-vsepr:resonance}{resonantiestructuur} $\mathrm{R_2C{=}O^+R}$ heeft
waarin elk atoom een octet heeft (\cref{prop:b1:electronic-effects:carbocation-order}).
Zuur wordt verbruikt bij de protoneringen en teruggevormd bij de
deprotoneringen: het is een \omterm{def:b1:elementary-steps:catalyst}{katalysator}.
\end{proof}

\begin{omfigure}
\setchemfig{atom sep=1.9em}
\schemestart
\chemfig{C(-[:150]R)(-[:210]H)=O}
\arrow{->[\ce{H+}][]}
\chemfig{C^{+}(-[:150]R)(-[:210]H)-OH}
\arrow{->[$\mathrm{R'OH}$][\ce{-H+}]}
\chemfig{C(-[:150]R)(-[:210]H)(-[2]OR')-OH}
\schemestop

\medskip
\schemestart
\chemfig{C(-[:150]R)(-[:210]H)(-[2]OR')-OH}
\arrow{->[\ce{H+}][\ce{-H2O}]}
\chemfig{C^{+}(-[:150]R)(-[:210]H)-OR'}
\arrow{->[$\mathrm{R'OH}$][\ce{-H+}]}
\chemfig{C(-[:150]R)(-[:210]H)(-[2]OR')-OR'}
\schemestop
\par\vspace{6pt}

{\small Acetaalvorming, in twee regels: eerst het \omterm{def:b1:acetals-protection:hemiacetal}{hemiacetaal} (activering,
additie van de alcohol, verlies van een proton), daarna het \omterm{def:b1:acetals-protection:hemiacetal}{acetaal}
(protonering van de OH, verlies van water tot het gestabiliseerde kation,
additie van een tweede alcohol, verlies van een proton).}
\end{omfigure}

De totale reactie, $\mathrm{RCHO} + 2\,\mathrm{R'OH} \rightleftharpoons \mathrm{RCH(OR')_2} + \ce{H2O}$, heeft
een bescheiden evenwichtsconstante, en water is een product.

\begin{proposition}[De acetalisering voortdrijven]\label{prop:b1:acetals-protection:dean-stark}
Als het gevormde water continu verwijderd wordt, gaat de acetalisering
door tot de carbonylverbinding (of de alcohol) opgebruikt is.
\end{proposition}

\begin{proof}
Het \omterm{def:b1:extent-q-and-k:quotient}{reactiequotiënt} $Q = [\text{acetaal}][\ce{H2O}]/([\text{aldehyde}]
[\mathrm{R'OH}]^2)$ blijft onder $K$ zolang water verwijderd wordt:
volgens \cref{ch:b1:extent-q-and-k} blijft de reactie naar rechts
verlopen, wat de waarde van $K$ ook is. Met een diol (ethaan-1,2-diol) is
het \omterm{def:b1:acetals-protection:hemiacetal}{acetaal} cyclisch, en is het evenwicht bovendien gunstiger, omdat één
molecuul diol er twee van de alcohol vervangt.
\end{proof}

\begin{omfigure}
\begin{tikzpicture}[font=\small]
  \pic at (0,0) {heatingmantle};
  \pic at (0,0) {rbflask={0.4}{omDef!14}};
  \pic (d) at (0,3.0) {deanstark={0.35}{omDef!35}{orange!15}};
  \pic at (0,6.4) {condenser};
  \draw[->] (2.4,4.0) node[right, align=left] {opvanger: water (onderste laag)\\keert niet terug; tolueen\\vloeit terug naar de kolf} -- (0.8,3.8);
  \draw[->] (-2.2,1.0) node[left, align=right] {aldehyde, diol, tolueen,\\zure katalysator, kokend} -- (-0.7,0.9);
  \draw[->] (1.8,8.0) node[right] {koeler} -- (0.4,7.8);
\end{tikzpicture}

{\small Een Dean--Stark-opvanger. De dampen van tolueen en water
condenseren en vallen in de gegradueerde opvanger; water, dichter dan
tolueen en er niet mee \omterm{def:b1:intermolecular-forces-solvents:miscibility}{mengbaar}, zakt naar beneden en blijft daar, terwijl
tolueen terugvloeit naar de kolf. Het opgevangen volume water meet hoe
ver de reactie gevorderd is.}
\end{omfigure}

\begin{proposition}[Stabiliteit van acetalen]\label{prop:b1:acetals-protection:acetal-stability}
\omterm{def:b1:acetals-protection:hemiacetal}{Acetalen} zijn stabiel in neutraal en basisch milieu en tegenover
\omterm{def:b1:electronic-effects:nucleophile}{nucleofielen}, \omterm{def:b1:organomagnesium:organometallic}{grignardreagentia} en hydridereagentia; waterig zuur
hydrolyseert ze terug tot de carbonylverbinding en de alcohol.
\end{proposition}

\begin{proof}
Een acetaalkoolstofatoom draagt geen \omterm{def:b1:electronic-effects:nucleophile}{vertrekkende groep} die een base of
een \omterm{def:b1:electronic-effects:nucleophile}{nucleofiel} zou kunnen verdrijven: \ce{RO-}, een \omterm{def:b1:predominant-reaction:strong-acid}{sterke base},
vertrekt niet, en er is geen C=O meer om aan te vallen. In waterig zuur
loopt het mechanisme hierboven achterwaarts: de grote overmaat water
maakt $Q < K$ voor de omgekeerde reactie.
\end{proof}

\section{Beschermen en ontschermen}

\begin{definition}[Beschermende groep]\label{def:b1:acetals-protection:protecting-group}
Een \emph{beschermende groep}\index{beschermende groep} is een groep die
in een molecuul ingevoerd wordt om een functie tijdelijk te maskeren,
zodat ze niet reageert in een stap die op een ander deel van het molecuul
gericht is. Het invoeren ervan is de
\emph{bescherming}\index{bescherming}, het verwijderen de
\emph{ontscherming}\index{ontscherming}.
\end{definition}

\begin{method}[Beschermen, laten reageren, ontschermen]\label{met:b1:acetals-protection:protect}
\begin{enumerate}
  \item Zoek de functie die in de geplande stap zou reageren, of het
        reagens zou vernietigen.
  \item Kies een \omterm{def:b1:acetals-protection:protecting-group}{beschermende groep} die in die stap stabiel is en
        verwijderd kan worden onder omstandigheden die de rest van het
        molecuul verdraagt (voor een aldehyde of keton tegenover basen,
        \omterm{def:b1:electronic-effects:nucleophile}{nucleofielen} of hydriden: een cyclisch \omterm{def:b1:acetals-protection:hemiacetal}{acetaal}).
  \item Bescherm; voer de geplande stap uit; ontscherm.
  \item Reken de kosten: twee stappen meer, elk met zijn \omterm{def:b1:extent-q-and-k:fractional-extent}{rendement}.
\end{enumerate}
\end{method}

\begin{example}[Een keton dat in de weg zit]\label{ex:b1:acetals-protection:ketoester}
Om een \omterm{def:b1:organomagnesium:organometallic}{grignardreagens} te laten adderen aan de aldehydegroep van een
molecuul dat ook een keton draagt, bestaat er geen eenvoudige
acetaalkeuze (aldehyden vormen gemakkelijker \omterm{def:b1:acetals-protection:hemiacetal}{acetalen} dan ketonen, zodat
de verkeerde groep beschermd zou worden); de volgorde van de stappen of
de reagentia moeten veranderen. \omterm{def:b1:acetals-protection:protecting-group}{Bescherming} is een strategie, geen
reflex.
\end{example}

\section{Oefeningen}

\begin{exercise}[$\star$]\label{exo:b1:acetals-protection:1}
Herken de koolstofatomen van een \omterm{def:b1:acetals-protection:hemiacetal}{hemiacetaal}, \omterm{def:b1:acetals-protection:hemiacetal}{acetaal}, hydraat, ether of
alcohol in: \ce{CH3CH(OH)OCH3}, \ce{CH3CH(OCH3)2}, \ce{CH3CH(OH)2},
\ce{CH3OCH3}, \ce{(CH3)2C(OCH2CH3)2}.
\end{exercise}

\begin{exercise}[$\star$]\label{exo:b1:acetals-protection:2}
Schrijf de totale vergelijking op van de vorming van het \omterm{def:b1:acetals-protection:hemiacetal}{acetaal} van
propanal met methanol, en van het cyclische \omterm{def:b1:acetals-protection:hemiacetal}{acetaal} van propanon met
ethaan-1,2-diol.
\end{exercise}

\begin{exercise}[$\star$]\label{exo:b1:acetals-protection:3}
Welke van deze reagentia laten een \omterm{def:b1:acetals-protection:hemiacetal}{acetaal} onveranderd: natronloog,
methylmagnesiumbromide, verdund zwavelzuur in water,
natriumboorhydride?
\end{exercise}

\begin{exercise}[$\star$]\label{exo:b1:acetals-protection:4}
Teken het cyclische \omterm{def:b1:acetals-protection:hemiacetal}{hemiacetaal} dat 4-hydroxybutanal vormt. Hoe groot is
de ring?
\end{exercise}

\begin{exercise}[$\star\star$]\label{exo:b1:acetals-protection:5}
Schrijf het volledige mechanisme op van de zuurgekatalyseerde vorming van
het dimethylacetaal van ethanal, met \omterm{def:b1:electronic-effects:curly-arrow}{gebogen pijlen}, en toon aan dat het
zuur teruggevormd wordt.
\end{exercise}

\begin{exercise}[$\star\star$]\label{exo:b1:acetals-protection:6}
Schrijf het mechanisme op van de zure hydrolyse van
2,2-dimethoxypropaan. Waarom gebruikt men een grote hoeveelheid water?
\end{exercise}

\begin{exercise}[$\star\star$]\label{exo:b1:acetals-protection:7}
Een acetalisering van \qty{0.250}{mol} van een keton wordt gevolgd met
een Dean--Stark-opvanger. Welk volume water zou bij volledige omzetting
opgevangen moeten worden (dichtheid ongeveer \qty{1.0}{g/mL})? Na een uur
is er \qty{3.2}{mL} opgevangen: wat is de omzetting?
\end{exercise}

\begin{exercise}[$\star\star$]\label{exo:b1:acetals-protection:8}
Verklaar waarom een cyclisch \omterm{def:b1:acetals-protection:hemiacetal}{acetaal} met ethaan-1,2-diol zich
vollediger vormt dan het \omterm{def:b1:acetals-protection:hemiacetal}{acetaal} met twee moleculen methanol, voor
dezelfde carbonylverbinding.
\end{exercise}

\begin{exercise}[$\star\star$]\label{exo:b1:acetals-protection:9}
Waarom moet de zure \omterm{def:b1:elementary-steps:catalyst}{katalysator} verwijderd (of geneutraliseerd) worden
voordat een beschermde verbinding met een \omterm{def:b1:organomagnesium:organometallic}{grignardreagens} behandeld
wordt?
\end{exercise}

\begin{exercise}[$\star\star\star$]\label{exo:b1:acetals-protection:10}
Stel een synthese voor van 5-hydroxy-5-methylhexaan-2-on,
\ce{(CH3)2C(OH)CH2CH2COCH3}, uit 4-chloorbutaan-2-on en propanon, met een
\omterm{def:b1:acetals-protection:protecting-group}{bescherming}. Geef elke stap, zijn reagentia en de rol van elk ervan.
\end{exercise}

\begin{exercise}[$\star\star\star$]\label{exo:b1:acetals-protection:11}
Toon met het \omterm{def:b1:extent-q-and-k:quotient}{reactiequotiënt} aan dat een grote overmaat methanol de
acetalisering van een aldehyde ook voortdrijft, en vergelijk deze methode
met het verwijderen van water.
\end{exercise}

\begin{exercise}[$\star\star\star$]\label{exo:b1:acetals-protection:12}
Verklaar waarom het cyclische \omterm{def:b1:acetals-protection:hemiacetal}{hemiacetaal} van glucose in zure methanol
aangevallen wordt en een methylacetaal (een methylglycoside) geeft,
terwijl de andere OH-groepen van glucose niet in ethers omgezet worden.
\end{exercise}

\section{Probleem: een grignardreagens dat een aldehyde draagt}

\begin{problem}[{Weekendprobleem --- waarom 4-broombutanal geen
grignardreagens kan geven, zijn bescherming als cyclisch acetaal met een
Dean--Stark-opvanger, de grignardadditie aan propanon, de ontscherming,
en het opgevangen volume water}]
\label{pb:b1:acetals-protection:1}
Een chemica heeft 5-hydroxy-5-methylhexanal nodig,
\ce{(CH3)2C(OH)CH2CH2CH2CHO}, en wil het maken uit 4-broombutanal,
\ce{BrCH2CH2CH2CHO}, en propanon. Ze vertrekt van \qty{15.09}{g}
4-broombutanal. Molaire massa's (\unit{g/mol}): \ce{C4H7BrO} 150.9,
\ce{C2H6O2} 62.0, \ce{C6H11BrO2} 194.9, \ce{C3H6O} 58.0, \ce{C9H18O3}
174.0, \ce{C7H14O2} 130.0, \ce{H2O} 18.0; dichtheid van water
\qty{0.997}{g/mL}.
% ledger: rho:water

\medskip
\noindent\textbf{Deel I --- Het probleem.}
\begin{enumerate}
  \item Schrijf het \omterm{def:b1:organomagnesium:organometallic}{grignardreagens} op dat 4-broombutanal zou geven.
  \item Waarom kan dat reagens niet bestaan? Schrijf de reactie op die het
        vernietigt.
  \item Welke functie moet beschermd worden, en waartegen?
  \item Waarom is een \omterm{def:b1:acetals-protection:hemiacetal}{acetaal} geschikt?
  \item Waarom verkiest men een cyclisch \omterm{def:b1:acetals-protection:hemiacetal}{acetaal} (met ethaan-1,2-diol)?
  \item Bereken de hoeveelheid stof 4-broombutanal.
  \item Bereken de massa diol voor een overmaat van \qty{20}{\percent}.
\end{enumerate}

\medskip
\noindent\textbf{Deel II --- \omterm{def:b1:acetals-protection:protecting-group}{Bescherming}.}
\begin{enumerate}[resume]
  \item Schrijf de vergelijking op van de acetalisering met
        ethaan-1,2-diol.
  \item Wat is de rol van de zure \omterm{def:b1:elementary-steps:catalyst}{katalysator}
        (4-methylbenzeensulfonzuur)?
  \item Beschrijf de werking van de Dean--Stark-opvanger.
  \item Verklaar met $Q$ en $K$ waarom het verwijderen van water de
        reactie voortdrijft.
  \item Waarom wordt het tolueen, en niet het water, naar de kolf
        teruggevoerd?
  \item Hoe weet de chemica dat de reactie volledig is?
  \item Bereken de massa beschermd bromide die bij volledige omzetting
        verwacht wordt.
\end{enumerate}

\medskip
\noindent\textbf{Deel III --- De grignardstap.}
\begin{enumerate}[resume]
  \item Schrijf de vorming op van het \omterm{def:b1:organomagnesium:organometallic}{grignardreagens} uit het beschermde
        bromide.
  \item Schrijf zijn additie aan propanon op en het gevormde \omterm{def:b1:alcohol-activation:alkoxide}{alkoxide}.
  \item Waarom moet de opwerking van deze stap licht basisch of neutraal
        zijn (bijvoorbeeld een waterige oplossing van ammoniumchloride) in
        plaats van zuur?
  \item Bereken de massa beschermde alcohol bij een \omterm{def:b1:extent-q-and-k:fractional-extent}{rendement} van
        \qty{70}{\percent} voor deze stap.
\end{enumerate}

\medskip
\noindent\textbf{Deel IV --- \omterm{def:b1:acetals-protection:protecting-group}{Ontscherming}.}
\begin{enumerate}[resume]
  \item Schrijf de \omterm{def:b1:acetals-protection:protecting-group}{ontscherming} in waterig zuur op.
  \item Waarom overleeft de tertiaire alcohol mild waterig zuur?
  \item Het product kan een zesring sluiten door zijn OH aan het aldehyde
        te adderen. Teken dat cyclische \omterm{def:b1:acetals-protection:hemiacetal}{hemiacetaal}.
  \item Bereken de massa hydroxyaldehyde bij een \omterm{def:b1:extent-q-and-k:fractional-extent}{rendement} van
        \qty{90}{\percent} voor de \omterm{def:b1:acetals-protection:protecting-group}{ontscherming}.
  \item Welk totaalrendement vanaf 4-broombutanal bereikt de chemica?
  \item Bereken de massa water die bij de beschermingsstap bij volledige
        omzetting gevormd wordt.
  \item Geef het volume water dat bij volledige omzetting in de
        Dean--Stark-opvanger opgevangen wordt.
\end{enumerate}
\end{problem}
